\subsection{Fonctions et exigences du système}
Pour pallier le manque de documentation et gérer la complexité, trois stratégies de 
communication ont été identifiées. Elles sont classées par ordre de préférence 
et font partie intégrante des exigences techniques du système. De plus, l'implémentation physique 
se déclinera sur différents modules d'évaluation (EKH05 et EKH19) qui implémentent cette puce.
La station de réception (RX) s'appuiera probablement sur le modèle EKH19 afin de pouvoir exploiter le système 
d'exploitation OpenWrt. Quant au Transmetteur (TX), il sera probablement basé sur le modèle EKH05 piloté par 
FreeRTOS.\newline

Les stratégies de communication réseau listées ci-dessous constituent des hypothèses de travail définies en début de 
projet. Elles ne sont pas figées et servent de base d'investigation. L'objectif principal reste l'optimisation des 
performances, ce qui inclut la liberté de pivoter vers une solution hybride si les tests sur le matériel révèlent 
une approche plus efficace.

\begin{itemize}
    \item \textbf{Option A (Privilégiée) :} Protocole UDP Broadcast côté TX
    combiné à une réception en mode moniteur (Monitor Mode) côté RX.
    \item \textbf{Option B (Alternative performance) :} Diffusion brute des données (Raw Broadcast, pas de protocole) 
    avec LDPC. C'est l'option la plus optimisée mais la plus complexe techniquement.
    \item \textbf{Option C (Sécurité) :} Utilisation des exemples standards du constructeur. Solution techniquement simple mais avec des performances potentiellement trop faibles.
\end{itemize}

\vspace{0.5cm}

\begin{longtable}{|p{6cm}|p{9cm}|}
\caption{Fonctions de service et exigences du système à concevoir}
\label{tab:fonctions_service} \\
\hline
\textbf{Fonctions de service} & \textbf{Exigences} \\
\hline
\endfirsthead

\multicolumn{2}{c}%
{{\bfseries \tablename\ \thetable{} -- suite de la page précédente}} \\
\hline
\textbf{Fonctions de service} & \textbf{Exigences} \\
\hline
\endhead

\hline \multicolumn{2}{|r|}{{Suite à la page suivante}} \\ \hline
\endfoot

\hline
\endlastfoot

\textbf{FS 1 :} Transmettre des données sans fil (MVP) 
& \textbf{E 1 :} Utiliser la norme IEEE 802.11ah (Wi-Fi HaLow). \\
\cline{2-2}
& \textbf{E 2 :} Forcer l'utilisation d'un MCS bas (index 1 ou 2). \\
\cline{2-2}
& \textbf{E 3 :} Activer la correction d'erreur LDPC. \\
\hline
\textbf{FS 2 :} Transmettre un flux vidéo (Objectif Final) 
& \textbf{E 4 :} Interfacer la puce MM8108 avec la carte Jetson (l'encodeur vidéo). \\
\cline{2-2}
& \textbf{E 5 :} Assurer le décodage et l'affichage sur la station de réception. \\
\end{longtable}

\vspace{0.5cm}

\begin{longtable}{|p{6cm}|p{9cm}|}
\caption{Fonctions techniques et exigences du système à concevoir}
\label{tab:fonctions_techniques} \\
\hline
\textbf{Fonctions techniques} & \textbf{Exigences} \\
\hline
\endfirsthead

\multicolumn{2}{c}%
{{\bfseries \tablename\ \thetable{} -- suite de la page précédente}} \\
\hline
\textbf{Fonctions techniques} & \textbf{Exigences} \\
\hline
\endhead

\hline \multicolumn{2}{|r|}{{Suite à la page suivante}} \\ \hline
\endfoot

\hline
\endlastfoot

\textbf{FT 1 :} Assurer la communication 
& \textbf{E 6 :} Implémenter la communication selon l'Option A (UDP Broadcast / Monitor mode) en priorité. (Ou solution permettant des performances similaires) \\
\cline{2-2}
& \textbf{E 7 :} Les options B (Raw) et C (Exemples) serviront de plans secondaires. \\
\hline
\textbf{FT 2 :} S'adapter aux limitations matérielles 
& \textbf{E 8 :} L'architecture logicielle doit être compatible avec les différentes variantes de modules basés sur la puce MM8108 (EKH05 et EKH19).\\
\end{longtable}

\vspace{0.5cm}

\begin{longtable}{|p{6cm}|p{9cm}|}
\caption{Fonctions de contrainte et exigences du système à concevoir}
\label{tab:fonctions_contrainte} \\
\hline
\textbf{Fonctions de contrainte} & \textbf{Exigences} \\
\hline
\endfirsthead

\multicolumn{2}{c}%
{{\bfseries \tablename\ \thetable{} -- suite de la page précédente}} \\
\hline
\textbf{Fonctions de contrainte} & \textbf{Exigences} \\
\hline
\endhead

\hline \multicolumn{2}{|r|}{{Suite à la page suivante}} \\ \hline
\endfoot

\hline
\endlastfoot

\textbf{FC 1 :} Respecter la contrainte temps 
& \textbf{E 9 :} Réalisable en 450 heures. Le MVP fait foi si l'objectif final (idéal) n'a pas pu être réalisé. \\
\hline
\textbf{FC 2 :} Composer avec la documentation limitée 
& \textbf{E 10 :} Justifier les choix techniques basés sur l'ingénierie inverse et les tests empiriques. \\
\end{longtable}